\chapter{Annexes}

\section{Extrait du script de collecte (\texttt{src/scraper.py})}

\begin{lstlisting}[language=Python, caption={Extraction structurée d'une annonce à partir de ses attributs HTML \texttt{data-t-listing\_*}}]
def parse_listing_card(anchor) -> dict:
    data = anchor.attrs

    bedrooms_tag = anchor.select_one(
        ".listing-card__header__tags__item--no-of-bedrooms")
    surface_tag = anchor.select_one(
        ".listing-card__header__tags__item--square-metres")
    bedrooms = extract_int_from_class(
        bedrooms_tag.get("class", []),
        "listing-card__header__tags__item--no-of-bedrooms_"
    ) if bedrooms_tag else None
    surface = extract_int_from_class(
        surface_tag.get("class", []),
        "listing-card__header__tags__item--square-metres_"
    ) if surface_tag else None

    return {
        "listing_id": data.get("data-t-listing_id"),
        "title": data.get("data-t-listing_title"),
        "category": data.get("data-t-listing_category_title"),
        "price_fcfa": data.get("data-t-listing_price"),
        "currency": data.get("data-t-listing_currency"),
        "bedrooms": bedrooms,
        "surface_m2": surface,
        "url": data.get("href"),
    }
\end{lstlisting}

\section{Pipeline de prétraitement et d'entraînement (\texttt{scikit-learn})}

\begin{lstlisting}[language=Python, caption={Pipeline de prétraitement commun à tous les modèles}]
preprocessor = ColumnTransformer(
    transformers=[
        ("num", StandardScaler(), NUMERIC_FEATURES),
        ("cat", OneHotEncoder(handle_unknown="ignore"),
         CATEGORICAL_FEATURES),
    ]
)

rf_pipeline = Pipeline([
    ("preprocessor", preprocessor),
    ("model", RandomForestRegressor(random_state=RANDOM_STATE)),
])

rf_search = RandomizedSearchCV(
    rf_pipeline,
    param_distributions=rf_param_distributions,
    n_iter=30,
    cv=KFold(5, shuffle=True, random_state=RANDOM_STATE),
    scoring="neg_mean_absolute_error",
    random_state=RANDOM_STATE,
    n_jobs=-1,
)
\end{lstlisting}

\section{Liste des quartiers conservés individuellement}

Les 33 quartiers suivants comptent chacun au moins 5 annonces dans le jeu de données nettoyé et
sont conservés comme catégories individuelles (les quartiers moins représentés sont regroupés
sous la catégorie \texttt{Autre}) : Almadies, Point-E, Mermoz, Mamelles, Fann, Ngor, Virage,
Saly, Ngaparou, Sacré-Cœur, Amitié, Dakar, Plateau, Mbao, Nguering, VDN, Ouakam, Somone, Lac
Rose, Yoff, et treize autres quartiers de moindre fréquence.

\section{Structure du dépôt de code}

\begin{lstlisting}
MemoireM1/
  src/scraper.py                 Collecte des annonces
  data/raw/                      Donnees brutes scrapees
  data/processed/                Donnees nettoyees
  notebooks/
    01_eda_nettoyage.ipynb       Nettoyage + analyse exploratoire
    02_modelisation.ipynb        Modelisation + optimisation
  models/                        Modeles entraines + metadonnees
  app.py                         Interface Streamlit
  requirements.txt               Dependances
  rapport/                       Present document (LaTeX)
\end{lstlisting}
